% ---------- ME242TB Materials Science for Engineers (CIE - Test 1, 02/04/2025)
\begin{xpaper}
\paperinfo{deptline={Department of Mechanical Engineering},date={02/04/2025},exam={CIE -- Test 1},maxmarks={10 + 50},sem={IV},prog={UG},duration={120 Min},
 title={Material Science for Engineers},code={ME242TB}}
\xhead
\begin{xtable}{Questions}
\qpart{Part -- A}
\q{1}{According to which principle the electrons fill the lowest energy orbital first before moving to higher energy ones.}{1}{1}{1}
\q{2}{In \blank[1.8cm] material the conduction and valence bands overlap so electrons can easily pass into the conduction band.}{1}{1}{1}
\q{3}{Define unit cell.}{1}{1}{1}
\q{4}{Write the classification of composites based on matrix.}{2}{1}{1}
\q{5}{The Burgers Vector is parallel to the dislocation line, the type of line defect is \blank[1.8cm].}{1}{1}{1}
\q{6}{If an atom is missing from its normal site in the matrix, the defect is called \blank[1.8cm] defect.}{1}{1}{2}
\q{7}{Define thermal conductivity with equation.}{2}{2}{2}
\q{8}{In dielectric materials the movement of electrons in response to an external electric field is \blank[1.8cm] type of polarization mechanism.}{1}{1}{2}
\qpart{Part -- B}
\q{1(a)}{With the help of neat sketches explain the Primary Interatomic bonds with examples.}{10}{2}{1}
\q{2(a)}{Classify solid materials explain them briefly with examples.}{10}{2}{1}
\q{3(a)}{Calculate Atomic packing factor for FCC unit cell.}{5}{3}{2}
\q{3(b)}{Classify and explain different crystal imperfections.}{5}{2}{2}
\q{4(a)}{Compare the properties and applications of metals and polymers.}{10}{3}{1}
\q{5(a)}{Illustrate the following thermoelectric effects: i) Seebeck effect ii) Peltier effect.}{5}{3}{2}
\q{5(b)}{Discuss different types of Insulating materials highlighting their applications.}{5}{2}{2}
\end{xtable}
\end{xpaper}
